\begin{python0}
from solutions import *; clear()
\end{python0}
\sectionthree{Blocks}

Of course it's not too surprising that you can have \EMPHASIZE{blocks} with the \texttt{for}-loop:
\begin{console}
std::cout << "entering the for-loop ...\n";
for (int i = 10; i > 0; i--)
{
        std::cout << "entered the block ...\n";
        std::cout << "i: " << i;
        std::cout << "exiting the block ...\n";
}
std::cout << "done with the for-loop ...\n";
\end{console}

It shouldn't be surprising that you can do this:
\begin{console}
std::cout << "entering the for-loop ...\n";
int x = 10, y = 0;
for (int i = x; i > y; i--)
{
        std::cout << "entered the block ...\n";
        std::cout << "i: " << i;
        std::cout << "exiting the block ...\n";
}
std::cout << "done with the for-loop ...\n";
\end{console}

By the way, just like in the case of \texttt{if-} and \texttt{if-else} statements, C/C++ programming tend to use blocks for the for-loop even when the body of the for-loop has only one statement:
\begin{console}
for (int i = 0; i < 10; i++)
{
        std::cout << i << std::endl;
}
\end{console}
\begin{ex}
 Remember our program that prints a table of squares?
\begin{console}
n   n^2
--- ---
0   0
1   1
2   4
\end{console}
Write a program that prints the squares of all positive integers from 0
to 100. Your program must be less than 10 lines long and you have 3
minutes to do that - so I don't mean you have a very
long \texttt{std::cout}!!! Here's some pseudocode that
might help:
\begin{console}
Print "n   n^2\n"
Print "--- ---\n"
For n running from 0 to 100 (inclusive):
    print n, spaces, square of n, and newline
\end{console}
\end{ex}

\begin{ex}
Modify your program so that it prompts the user for integers
and assign them to integer variables \texttt{start} and \texttt{end}. Print
a table of squares from \texttt{start} to \texttt{end}.
\end{ex}
\begin{ex}
Can you use a \texttt{double} variable to control the
for-loop? Try this
\begin{console}
std::cout << "before" << std::endl;

for (double x = 3.1; x < 7.2; x += 0.1)
{
        std::cout << x << std::endl;
}

std::cout << "after" << std::endl;
\end{console}
\end{ex}

\begin{ex}
Write a program that prints a table of cubes (like
the program on squares above).
\end{ex}

